% Geographic Stability
\label{sec:stability}

\subsection{Boundary Overlap Analysis}

We compute Intersection-over-Union (IoU) for districts across consecutive census cycles. Table~\ref{tab:stability} summarizes boundary persistence:

\begin{table}[h]
\centering
\begin{tabular}{lrr}
\toprule
\textbf{Transition} & \textbf{Mean IoU} & \textbf{\% Districts IoU $>$ 0.7} \\
\midrule
2000 $\rightarrow$ 2010 & 0.68 & 54\% \\
2010 $\rightarrow$ 2020 & 0.71 & 61\% \\
\bottomrule
\end{tabular}
\caption{District boundary persistence across census cycles}
\label{tab:stability}
\end{table}

\textbf{Key finding}: Geographic stability increased slightly from 2000--2010 to 2010--2020, despite significant seat reallocations. This suggests that algorithmic redistricting can adapt to population shifts while maintaining substantial boundary continuity.

An IoU threshold of 0.7 indicates that 70\% of a district's area remains unchanged. By this criterion, 61\% of districts in 2020 were substantially similar to their 2010 counterparts—despite 13 states gaining or losing seats.

\subsection{Population Reassignment}

We measure the fraction of population reassigned to different district numbers:

\begin{itemize}
    \item 2000 $\rightarrow$ 2010: 28\% of population reassigned
    \item 2010 $\rightarrow$ 2020: 24\% of population reassigned
\end{itemize}

\textbf{Interpretation}: Algorithmic redistricting maintains structural continuity. Most residents (76\% in 2010--2020) remain in districts geographically similar to their prior representation. This is notable given that population growth is geographically uneven—some regions grow rapidly while others stagnate.

\subsection{State-Level Variability}

Figure~\ref{fig:stability-map} shows state-level IoU scores as a choropleth map. High-stability states (IoU $>$ 0.75) include:

\begin{itemize}
    \item \textbf{Single-district states}: VT, WY, ND, SD, MT, DE (boundaries identical by definition)
    \item \textbf{Slow-growth states}: RI, WV (minimal population shifts)
    \item \textbf{Geographically constrained}: HI, AK (islands/terrain limit boundary options)
\end{itemize}

Low-stability states (IoU $<$ 0.50) include:

\begin{itemize}
    \item \textbf{High-growth states}: TX, FL, AZ (population shifts force major redrawing)
    \item \textbf{Seat-change states}: NY, PA, OH (losing seats requires merging districts)
\end{itemize}

Texas exemplifies the stability-responsiveness trade-off: IoU = 0.42 (low stability) reflects gaining 6 seats and absorbing 4.7M new residents. The algorithm preserves districts in stable regions while redrawing high-growth areas.

\subsection{Temporal Trends}

The modest increase in stability from 2000--2010 (IoU = 0.68) to 2010--2020 (IoU = 0.71) likely reflects:

\begin{enumerate}
    \item \textbf{Demographic deceleration}: U.S. population growth slowed from 13.2\% (1990--2000) to 9.7\% (2000--2010) to 7.4\% (2010--2020)
    \item \textbf{Smaller reallocations}: Fewer seats changed hands in 2020 (13 states) than 2010 (18 states)
    \item \textbf{Geographic concentration}: Growth increasingly concentrated in specific metros (Austin, Phoenix) rather than statewide
\end{enumerate}

\subsection{Comparison to Enacted Districts}

Enacted districts show lower stability (mean IoU = 0.51 for 2010--2020) than algorithmic districts (IoU = 0.71). This reflects partisan and incumbent-protection incentives that favor wholesale redrawing over incremental adjustment.

\textbf{Implication}: Algorithmic redistricting provides stability benefits beyond compactness—it adapts to demographic change without unnecessary disruption.
